%
% $Header: /home/cvs/root/haskell-report/report/modules.verb,v 1.20 2003/01/13 13:08:56 simonpj Exp $
%
%**<title>The Haskell 98 Report: Modules</title>
%*section 5
%**~header
\chapter{Modules} 
\label{modules} 
\index{module}

A module defines a collection of values, datatypes, type synonyms,
classes, etc.~(see Chapter~\ref{declarations}), in an environment created
by a set of {\em imports} (resources brought into scope from other modules).
It {\em exports} some of these resources, making them available to
other modules.  
We use the term {\em entity}\index{entity} to refer to
a value, type, or class defined in, imported into, or perhaps
exported from a module.

A \Haskell{} {\em program} is a collection of modules, one of
which, by convention, must be called \mbox{\tt Main}\indexmodule{Main} and must
export the value \mbox{\tt main}\indextt{main}.  The {\em value} of the program
is the value of the identifier \mbox{\tt main} in module \mbox{\tt Main},
which must be a computation of type $\mbox{\tt IO}~\tau$ for some type $\tau$
(see Chapter~\ref{io}).  When the program is executed, the computation
\mbox{\tt main} is performed, and its result (of type $\tau$) is discarded.

Modules may reference other modules via explicit
\mbox{\tt import} declarations, each giving the name of a module to be
imported and specifying its entities to be imported.
Modules may be mutually recursive.

Modules are used for name-space control, and are not first class values.
A multi-module Haskell program can be converted into a single-module
program by giving each entity a unique name, changing all occurrences
to refer to the appropriate unique name, and then concatenating all the module
bodies\footnote{There are two minor exceptions to this statement.
First, \mbox{\tt default} declarations scope over a single module (Section~\ref{default-decls}).
Second, Rule 2 of the monomorphism restriction (Section~\ref{sect:monomorphism-restriction})
is affected by module boundaries.
}.  
For example, here is a three-module program:
\bprog
{\small\begin{verbatim}
  module Main where
    import A
    import B
    main = A.f >> B.f

  module A where
    f = ...

  module B where
    f = ...
\end{verbatim}}
\eprog
It is equivalent to the following single-module program:
\bprog
{\small\begin{verbatim}
  module Main where
    main = af >> bf

    af = ...

    bf = ...
\end{verbatim}}
\eprog
Because they are allowed to be mutually recursive,
modules allow a program to be partitioned freely without regard to
dependencies.

\begin{haskellprime}
A module name (lexeme \mbox{$\it modid$}) is a sequence of one or more
identifiers beginning with capital letters, separated by dots, with no
intervening spaces.  For example, \mbox{\tt Data.Bool}, \mbox{\tt Main} and
\mbox{\tt Foreign.Marshal.Alloc} are all valid module names.

\begin{flushleft}\it\begin{tabbing}
\hspace{0.6in}\=\hspace{3.1in}\=\kill
$\it modid$\>\makebox[3.5em]{$\rightarrow$}$\it \hprime{\{conid\ \makebox{\tt .}\}}\ conid$\>\makebox[3em]{}$\it (\tr{modules})$
\end{tabbing}\end{flushleft}
\indexsyn{modid}%

Module names can be thought of as being arranged in a hierarchy in
which appending a new component creates a child of the original module
name.  For example, the module \mbox{\tt Control.Monad.ST} is a child of the
\mbox{\tt Control.Monad} sub-hierarchy.  This is purely a convention, however,
and not part of the language definition; in this report a \mbox{$\it modid$} is
treated as a single identifier occupying a flat namespace.
\end{haskellprime}

There is one distinguished module, \mbox{\tt Prelude}, which is imported into
all modules by default (see Section~\ref{standard-prelude}), plus a
set of standard library modules that may be imported as required
(see Part~\ref{libraries}).

\section{Module Structure} 
\label{module-implementations}

A module defines a mutually
recursive scope containing declarations for value bindings, data
types, type synonyms, classes, etc. (see Chapter~\ref{declarations}).

\begin{flushleft}\it\begin{tabbing}
\hspace{0.6in}\=\hspace{3.1in}\=\kill
$\it module$\>\makebox[3.5em]{$\rightarrow$}$\it \makebox{\tt module}\ modid\ [exports]\ \makebox{\tt where}\ body$\\ 
$\it $\>\makebox[3.5em]{$|$}$\it body$\\ 
$\it body$\>\makebox[3.5em]{$\rightarrow$}$\it \makebox{\tt {\char'173}}\ impdecls\ \makebox{\tt ;}\ topdecls\ \makebox{\tt {\char'175}}$\\ 
$\it $\>\makebox[3.5em]{$|$}$\it \makebox{\tt {\char'173}}\ impdecls\ \makebox{\tt {\char'175}}$\\ 
$\it $\>\makebox[3.5em]{$|$}$\it \makebox{\tt {\char'173}}\ topdecls\ \makebox{\tt {\char'175}}$\\ 
$\it $\\ 
$\it impdecls$\>\makebox[3.5em]{$\rightarrow$}$\it impdecl_1\ \makebox{\tt ;}\ \ldots \ \makebox{\tt ;}\ impdecl_n$\>\makebox[3em]{}$\it (n\geq 1)$\\ 
$\it topdecls$\>\makebox[3.5em]{$\rightarrow$}$\it topdecl_1\ \makebox{\tt ;}\ \ldots \ \makebox{\tt ;}\ topdecl_n$\>\makebox[3em]{}$\it (n\geq 1)$
\end{tabbing}\end{flushleft}

\indexsyn{module}%
\indexsyn{body}%
\indexsyn{impdecls}%
\indexsyn{topdecls}%

A module 
begins with a header: the keyword
\mbox{\tt module}, the module name, and a list of entities (enclosed in round
parentheses) to be exported.  The header is followed by a possibly-empty
list of \mbox{\tt import} declarations (\mbox{$\it impdecls$}, Section~\ref{import}) that specify modules to be imported,
optionally restricting the imported bindings.  
This is followed by a possibly-empty list of top-level declarations (\mbox{$\it topdecls$}, Chapter~\ref{declarations}).

An abbreviated form of module, consisting only 
of\index{abbreviated module}
the module body, is permitted.  If this is used, the header is assumed to be
`\mbox{\tt module\ Main(main)\ where}'.
If the first lexeme in the
abbreviated module is not a \mbox{\tt {\char'173}}, then the layout rule applies
for the top level of the module.

\section{Export Lists}
\label{export}
\index{export list}

\begin{flushleft}\it\begin{tabbing}
\hspace{0.6in}\=\hspace{3.1in}\=\kill
$\it exports$\>\makebox[3.5em]{$\rightarrow$}$\it \makebox{\tt (}\ export_1\ \makebox{\tt ,}\ \ldots \ \makebox{\tt ,}\ export_n\ [\ \makebox{\tt ,}\ ]\ \makebox{\tt )}$\>\makebox[3em]{}$\it (n\geq 0)$\\ 
$\it $\\ 
$\it export$\>\makebox[3.5em]{$\rightarrow$}$\it qvar$\\ 
$\it $\>\makebox[3.5em]{$|$}$\it qtycon\ [\makebox{\tt (..)}\ |\ \makebox{\tt (}\ cname_1\ \makebox{\tt ,}\ \ldots \ \makebox{\tt ,}\ cname_n\ \makebox{\tt )}]$\>\makebox[3em]{}$\it (n\geq 0)$\\ 
$\it $\>\makebox[3.5em]{$|$}$\it qtycls\ [\makebox{\tt (..)}\ |\ \makebox{\tt (}\ var_1\ \makebox{\tt ,}\ \ldots \ \makebox{\tt ,}\ var_n\ \makebox{\tt )}]$\>\makebox[3em]{}$\it (n\geq 0)$\\ 
$\it $\>\makebox[3.5em]{$|$}$\it \makebox{\tt module}\ modid$\\ 
$\it $\\ 
$\it cname$\>\makebox[3.5em]{$\rightarrow$}$\it var\ |\ con$
\end{tabbing}\end{flushleft}
\indexsyn{exports}%
\indexsyn{export}%

An {\em export list} identifies the entities to be exported by a
module declaration.  A module implementation may only export an entity
that it declares, or that it imports from some other module.  If the
export list is omitted, all values, types and classes defined in the
module are exported, {\em but not those that are imported}.

Entities in an export list may be named as follows:
\begin{enumerate}
\item
A value, field name, or class method, whether declared in
the module body or imported,
may be named by giving the name of the value as a \mbox{$\it qvarid$}, which must be in scope.
Operators should be enclosed in parentheses to turn them into
\mbox{$\it qvarid$}s.  

\item
An algebraic datatype \mbox{$\it T$}
\index{algebraic datatype}
declared by a \mbox{\tt data} or \mbox{\tt newtype} declaration may be named in one of
three ways: 
\begin{itemize}
\item
The form \mbox{$\it T$} names the type {\em but not the constructors or field names}.
The ability to export a type without its constructors allows the
construction of abstract datatypes (see Section~\ref{abstract-types}).
\item
The form $T\mbox{\tt (}c_1\mbox{\tt ,}\ldots\mbox{\tt ,}c_n\mbox{\tt )}$, 
names the type and some or all of its constructors and field names.  
% Restriction removed March 02:
% The subordinate names $c_i$ must not contain duplicates.  
\item
The abbreviated form \mbox{$\it T\makebox{\tt (..)}$} names the type 
and all its constructors and field names that are currently in scope
(whether qualified or not).
\end{itemize}
In all cases, the (possibly-qualified) type constructor \mbox{$\it T$} must be in scope. 
The constructor and field names $c_i$ in the second form are unqualified;
one of these subordinate names is legal if and only if (a) it names a constructor
or field of \mbox{$\it T$}, and (b) the constructor or field
is in scope in the module body {\em regardless of whether it is in scope
under a qualified or unqualified name}. For example, the following is 
legal
\bprog
{\small\begin{verbatim}
  module A( Mb.Maybe( Nothing, Just ) ) where
    import qualified Data.Maybe as Mb
\end{verbatim}}
\eprog
Data constructors cannot be named in export lists except as subordinate names, because
they cannot otherwise be distinguished from type constructors.

\item
A type synonym \mbox{$\it T$} declared by a
\mbox{\tt type} declaration may be named by the form \mbox{$\it T$}, where \mbox{$\it T$} is in scope.
\index{type synonym}

\item
\index{class declaration}
A class $C$ with operations $f_1,\ldots,f_n$
declared in a \mbox{\tt class} declaration may be named in one of three ways:
\begin{itemize}
\item
The form \mbox{$\it C$} names the class {\em but not the class methods}.
\item
The form $C\mbox{\tt (}f_1\mbox{\tt ,}\ldots\mbox{\tt ,}f_n\mbox{\tt )}$, names the class and some or all
of its methods.  
% Restriction removed March 02:
% The subordinate names $f_i$ must not contain duplicates.
\item
The abbreviated form $C\mbox{\tt (..)}$ names the class and all its methods
that are in scope (whether qualified or not).
\end{itemize}
In all cases, \mbox{$\it C$} must be in scope.  In the second form,
one of the (unqualified) subordinate names $f_i$ is legal if and only if (a) it names a
class method of \mbox{$\it C$}, and (b) the class method 
is in scope in the module body regardless of whether it is in scope
under a qualified or unqualified name.

\item
The form ``\mbox{\tt module\ M}'' names the set of all entities that are in
scope with both an unqualified name ``\mbox{\tt e}'' and a qualified name
``\mbox{\tt M.e}''.
This set may be empty.
For example:
\bprog
{\small\begin{verbatim}
  module Queue( module Stack, enqueue, dequeue ) where
      import Stack
      ...
\end{verbatim}}
\eprog
Here the module \mbox{\tt Queue} uses the module name \mbox{\tt Stack} in its export
list to abbreviate all the entities imported from \mbox{\tt Stack}.  

A module can name its own local definitions in its export
list using its own name in the ``\mbox{\tt module\ M}'' syntax, because a local
declaration brings into scope both a qualified and unqualified name (Section~\ref{qualifiers}). 
For example:
\bprog
{\small\begin{verbatim}
  module Mod1( module Mod1, module Mod2 ) where
  import Mod2
  import Mod3
\end{verbatim}}
\eprog
Here module \mbox{\tt Mod1} exports all local definitions as well as those
imported from \mbox{\tt Mod2} but not those imported from \mbox{\tt Mod3}.

It is an error to use \mbox{\tt module\ M} in an export list unless \mbox{\tt M} is 
the module bearing the export list, or \mbox{\tt M} is imported by at 
least one import declaration (qualified or unqualified).
\end{enumerate}
Exports lists are cumulative: the set of entities exported by an export
list is the union of the entities exported by the individual items of the list.

It makes no difference to an importing module how an entity was 
exported.  For example, a field name \mbox{\tt f} from data type \mbox{\tt T} may be exported individually
(\mbox{\tt f}, item (1) above); or as an explicitly-named member of its data type (\mbox{\tt T(f)}, item (2));
or as an implicitly-named member (\mbox{\tt T(..)}, item(2)); or by exporting an entire
module (\mbox{\tt module\ M}, item (5)).  

The {\em unqualified} names of the entities exported by a module must all be distinct
(within their respective namespace).  For example
\bprog
{\small\begin{verbatim}
  module A ( C.f, C.g, g, module B ) where   -- an invalid module
  import B(f)
  import qualified C(f,g)
  g = f True
\end{verbatim}}
\eprog
There are no name clashes within module \mbox{\tt A} itself, 
but there are name clashes in the export list between \mbox{\tt C.g} and \mbox{\tt g}
(assuming \mbox{\tt C.g} and \mbox{\tt g} are different entities -- remember, modules
can import each other recursively), and between \mbox{\tt module\ B} and \mbox{\tt C.f}
(assuming \mbox{\tt B.f} and \mbox{\tt C.f} are different entities).

\section{Import Declarations}
\label{import}
\index{import declaration}

\begin{flushleft}\it\begin{tabbing}
\hspace{0.6in}\=\hspace{3.1in}\=\kill
$\it impdecl$\>\makebox[3.5em]{$\rightarrow$}$\it \makebox{\tt import}\ [\makebox{\tt qualified}]\ modid\ [\makebox{\tt as}\ modid]\ [impspec]$\\ 
$\it $\>\makebox[3.5em]{$|$}$\it $\>\makebox[3em]{}$\it (empty\ declaration)$\\ 
$\it impspec$\>\makebox[3.5em]{$\rightarrow$}$\it \makebox{\tt (}\ import_1\ \makebox{\tt ,}\ \ldots \ \makebox{\tt ,}\ import_n\ [\ \makebox{\tt ,}\ ]\ \makebox{\tt )}$\>\makebox[3em]{}$\it (n\geq 0)$\\ 
$\it $\>\makebox[3.5em]{$|$}$\it \makebox{\tt hiding}\ \makebox{\tt (}\ import_1\ \makebox{\tt ,}\ \ldots \ \makebox{\tt ,}\ import_n\ [\ \makebox{\tt ,}\ ]\ \makebox{\tt )}$\>\makebox[3em]{}$\it (n\geq 0)$\\ 
$\it $\\ 
$\it import$\>\makebox[3.5em]{$\rightarrow$}$\it var$\\ 
$\it $\>\makebox[3.5em]{$|$}$\it tycon\ [\ \makebox{\tt (..)}\ |\ \makebox{\tt (}\ cname_1\ \makebox{\tt ,}\ \ldots \ \makebox{\tt ,}\ cname_n\ \makebox{\tt )}]$\>\makebox[3em]{}$\it (n\geq 0)$\\ 
$\it $\>\makebox[3.5em]{$|$}$\it tycls\ [\makebox{\tt (..)}\ |\ \makebox{\tt (}\ var_1\ \makebox{\tt ,}\ \ldots \ \makebox{\tt ,}\ var_n\ \makebox{\tt )}]$\>\makebox[3em]{}$\it (n\geq 0)$\\ 
$\it cname$\>\makebox[3.5em]{$\rightarrow$}$\it var\ |\ con$
\end{tabbing}\end{flushleft}
%             var
%          |  tycon
%          |  tycon \mbox{\tt (..)}
%          |  tycon \mbox{\tt (} con_1 \mbox{\tt ,} ... \mbox{\tt ,} con_n\mbox{\tt )} & (n>=1)
%          |  tycls \mbox{\tt (..)}
%          |  tycls \mbox{\tt (} var_1 \mbox{\tt ,} ... \mbox{\tt ,} var_n\mbox{\tt )} & (n>=0)
\indexsyn{impdecl}%
\indexsyn{impspec}%
\indexsyn{import}%
\indexsyn{cname}%

The entities exported by a module may be brought into scope in
another module with
an \mbox{\tt import}
declaration at the beginning
of the module.  
The \mbox{\tt import} declaration names the module to be
imported
and optionally specifies the entities to be imported.
A single module may be imported by more than one \mbox{\tt import} declaration.  
Imported names serve as top level declarations: they scope over
the entire body of the module but may  be shadowed by local
non-top-level bindings.  

The effect of multiple \mbox{\tt import} declarations is strictly
cumulative: an entity is in scope if it is imported by any of the \mbox{\tt import}
declarations in a module.  The ordering of import declarations is irrelevant.

Lexically, the terminal symbols ``\mbox{\tt as}'', ``\mbox{\tt qualified}'' and
``\mbox{\tt hiding}'' are each a \mbox{$\it varid$} rather than a \mbox{$\it reservedid$}.  They have
special significance only in the context of an \mbox{\tt import} declaration;
they may also be used as variables.

\subsection{What is imported}
\label{whatisimported}

Exactly which entities are to be imported can be specified in one
of the following three ways:\nopagebreak[4]
\begin{enumerate}
\item
The imported entities can be specified explicitly
by listing them in parentheses.
Items in the list have the same form as those in export lists, except
qualifiers are not permitted and
the `\mbox{\tt module} \mbox{$\it modid$}' entity is not permitted.  When the \mbox{\tt (..)} form
of import is used for a type or class, the \mbox{\tt (..)} refers to all of the
constructors, methods, or field names exported from the module.

The list must name only
entities exported by the imported module.
The list may be empty, in which case nothing except the instances is
imported.

\item
Entities can be excluded by 
using the form \mbox{\tt hiding(}\mbox{$\it import_1\ \makebox{\tt ,}\ \ldots \ \makebox{\tt ,}\ import_n$}
\mbox{\tt )},\index{hiding} which
specifies that all entities exported by the named module should
be imported except for those named in the list.  Data constructors may be
named directly in hiding lists without being prefixed by the
associated type.  Thus, in
\bprog
{\small\begin{verbatim}
  import M hiding (C)
\end{verbatim}}
\eprog
any constructor, class, or type named \mbox{\tt C} is excluded.  In contrast,
using \mbox{\tt C} in an import list names only a class or type.  

It is an error to hide an entity that is not, in fact, exported by
the imported module.

\item
Finally, if \mbox{$\it impspec$} is omitted then 
all the entities exported by the specified module are imported.
\end{enumerate}

\subsection{Qualified import}
\index{qualified name}

For each entity imported under the rules of Section~\ref{whatisimported},
the top-level environment is extended.  If the import declaration used
the \mbox{\tt qualified} keyword, only the {\em qualified name} of the entity is
brought into scope.  If the \mbox{\tt qualified} keyword is omitted, then {\em both} the
qualified {\em and} unqualified name of the entity is brought into scope.
Section~\ref{qualifiers} describes qualified names in more detail.

The qualifier on the imported name is either the name of the imported module,
or the local alias given in the \mbox{\tt as} clause (Section~\ref{as-clause}) 
on the \mbox{\tt import} statement.
Hence, {\em the qualifier is not necessarily the name of the module in which the
entity was originally declared}.

The ability to exclude the unqualified names allows full programmer control of
the unqualified namespace: a locally defined entity can share the same
name as a qualified import:
\par
{\small
\bprog
{\small\begin{verbatim}
  module Ring where
  import qualified Prelude    -- All Prelude names must be qualified
  import Data.List( nub )

  l1 + l2 = l1 Prelude.++ l2  -- This + differs from the one in the Prelude
  l1 * l2 = nub (l1 + l2)     -- This * differs from the one in the Prelude

  succ = (Prelude.+ 1)
\end{verbatim}}
\eprog
}

\subsection{Local aliases}
\label{as-clause}

Imported modules may be assigned a local alias in the importing module
using the \mbox{\tt as} clause.
For example, in
\bprog
{\small\begin{verbatim}
  import qualified VeryLongModuleName as C
\end{verbatim}}
\eprog
entities must be referenced using `\mbox{\tt C.}' as a qualifier instead of
`\mbox{\tt VeryLongModuleName.}'.  This also allows a different module to be substituted
for \mbox{\tt VeryLongModuleName} without changing the qualifiers used for the imported module.
It is legal for more than one module in scope to 
use the same qualifier, provided that all names can still be resolved unambiguously.
For example:
\bprog
{\small\begin{verbatim}
  module M where
    import qualified Foo as A
    import qualified Baz as A
    x = A.f
\end{verbatim}}
\eprog
This module is legal provided only that \mbox{\tt Foo} and \mbox{\tt Baz} do not both export \mbox{\tt f}.

An \mbox{\tt as} clause may also be used on an un-\mbox{\tt qualified} \mbox{\tt import} statement:
\bprog
{\small\begin{verbatim}
  import Foo as A(f)
\end{verbatim}}
\eprog
This declaration brings into scope \mbox{\tt f} and \mbox{\tt A.f}.


\subsection{Examples}

To clarify the above import rules, suppose the module \mbox{\tt A} exports \mbox{\tt x} and \mbox{\tt y}.
Then this table shows what names are brought into scope by the specified import statement:
\begin{center}
\begin{tabular}{|ll|}
\hline
Import declaration & Names brought into scope \\
\hline
  \mbox{\tt import\ A}                    & \mbox{\tt x}, \mbox{\tt y}, \mbox{\tt A.x}, \mbox{\tt A.y} \\
  \mbox{\tt import\ A()}                  & (nothing)     \\
  \mbox{\tt import\ A(x)}                 & \mbox{\tt x}, \mbox{\tt A.x} \\
  \mbox{\tt import\ qualified\ A}          & \mbox{\tt A.x}, \mbox{\tt A.y} \\
  \mbox{\tt import\ qualified\ A()}        & (nothing) \\
  \mbox{\tt import\ qualified\ A(x)}       & \mbox{\tt A.x} \\
  \mbox{\tt import\ A\ hiding\ ()}          & \mbox{\tt x}, \mbox{\tt y}, \mbox{\tt A.x}, \mbox{\tt A.y} \\
  \mbox{\tt import\ A\ hiding\ (x)}         & \mbox{\tt y}, \mbox{\tt A.y} \\
  \mbox{\tt import\ qualified\ A\ hiding\ ()}        & \mbox{\tt A.x}, \mbox{\tt A.y} \\
  \mbox{\tt import\ qualified\ A\ hiding\ (x)}       & \mbox{\tt A.y} \\
  \mbox{\tt import\ A\ as\ B}               & \mbox{\tt x}, \mbox{\tt y}, \mbox{\tt B.x}, \mbox{\tt B.y} \\
  \mbox{\tt import\ A\ as\ B(x)}            & \mbox{\tt x}, \mbox{\tt B.x} \\
  \mbox{\tt import\ qualified\ A\ as\ B}     & \mbox{\tt B.x}, \mbox{\tt B.y} \\
\hline
\end{tabular}
\end{center}
In all cases, all instance declarations in scope in module \mbox{\tt A} are imported
(Section~\ref{import-instances}).

\section{Importing and Exporting Instance Declarations}
\label{import-instances}
\index{instance declaration!importing and exporting}

Instance declarations cannot be explicitly named on import or export
lists.  All instances in scope within a module are {\em always}
exported and any import brings {\em all} instances in from the
imported module.  Thus, an
instance declaration is in scope if and only if a chain of \mbox{\tt import}
declarations leads to the module containing the instance declaration.

For example, \mbox{\tt import\ M()} does not bring
any new names in scope from module \mbox{\tt M}, but does bring in any instances
visible in \mbox{\tt M}.  A module whose only purpose is to provide instance
declarations can have an empty export list.  For example
\bprog
{\small\begin{verbatim}
  module MyInstances() where
    instance Show (a -> b) where
      show fn = "<<function>>"
    instance Show (IO a) where
      show io = "<<IO action>>"
\end{verbatim}}
\eprog

\section{Name Clashes and Closure}

\subsection{Qualified names}\index{qualified name}
\label{qualifiers}

A {\em qualified name} is written as \mbox{$\it modid$}\mbox{\tt .}\mbox{$\it name$} (Section~\ref{ids}).
A qualified name is brought into scope:
\begin{itemize}
\item {\em By a top level declaration.}
A top-level declaration brings into scope both the unqualified {\em and}
the qualified name of the entity being defined.  Thus:
\bprog
{\small\begin{verbatim}
  module M where
    f x = ...
    g x = M.f x x
\end{verbatim}}
\eprog
is legal.  The {\em defining} occurrence must mention the {\em unqualified} name; therefore, it is
illegal to write
\bprog
{\small\begin{verbatim}
  module M where
    M.f x = ...                 -- ILLEGAL
    g x = let M.y = x+1 in ...  -- ILLEGAL
\end{verbatim}}
\eprog
\item {\em By an \mbox{\tt import} declaration.}  An \mbox{\tt import} declaration, whether \mbox{\tt qualified} or not,
always brings into scope the qualified name of the imported entity (Section~\ref{import}).
This allows a qualified
import to be replaced with an unqualified one without forcing changes
in the references to the imported names.  
\end{itemize}


\subsection{Name clashes}

If a module contains a bound occurrence of a name, such as \mbox{\tt f} or \mbox{\tt A.f},
it must be possible unambiguously to resolve which entity is thereby referred to;
that is, there must be only one binding for \mbox{\tt f} or \mbox{\tt A.f} respectively.

It is {\em not} an error for there to exist names that cannot be so 
resolved, provided that the program does not mention those names.  For example:
\bprog
{\small\begin{verbatim}
  module A where
    import B
    import C
    tup = (b, c, d, x)
  
  module B( d, b, x, y ) where
    import D
    x = ...
    y = ...
    b = ...
  
  module C( d, c, x, y ) where
    import D
    x = ...
    y = ...
    c = ...

  module D( d ) where
    d = ...
\end{verbatim}}
\eprog
Consider the definition of \mbox{\tt tup}.  
\begin{itemize}
\item The references to \mbox{\tt b} and \mbox{\tt c}
can be unambiguously resolved to \mbox{\tt b} declared in \mbox{\tt B}, and \mbox{\tt c} declared in
\mbox{\tt C} respectively.  
\item The reference to \mbox{\tt d} is unambiguously resolved to
\mbox{\tt d} declared in \mbox{\tt D}.  In this case the same entity is brought into scope by two routes
(the import of \mbox{\tt B} and the import of \mbox{\tt C}), and can be referred to in \mbox{\tt A} by the names
\mbox{\tt d}, \mbox{\tt B.d}, and \mbox{\tt C.d}.
\item The reference to \mbox{\tt x} is ambiguous: it could mean \mbox{\tt x} declared in \mbox{\tt B}, or \mbox{\tt x} 
declared in \mbox{\tt C}.  The ambiguity could be fixed by replacing the reference to \mbox{\tt x} by
\mbox{\tt B.x} or \mbox{\tt C.x}.
\item There is no reference to \mbox{\tt y}, so it is not erroneous that distinct entities called
\mbox{\tt y} are exported by both \mbox{\tt B} and \mbox{\tt C}.  An error is only reported if \mbox{\tt y} is actually mentioned.
\end{itemize}

The name occurring in a type signature or fixity declarations is
always unqualified, and unambiguously refers to another declaration in
the same declaration list (except that the fixity declaration for a
class method can occur at top level --- Section~\ref{fixity}). For example,
the following module is legal:
\bprog
{\small\begin{verbatim}
  module F where

    sin :: Float -> Float
    sin x = (x::Float)

    f x = Prelude.sin (F.sin x)
\end{verbatim}}
\eprog
The local declaration for \mbox{\tt sin} is
legal, even though the Prelude function \mbox{\tt sin} is implicitly in
scope. The references to \mbox{\tt Prelude.sin} and \mbox{\tt F.sin} must both be qualified
to make it unambiguous which \mbox{\tt sin} is meant. However, the unqualified
name \mbox{$\it \makebox{\tt sin}$} in the type signature in the first line of \mbox{\tt F} unambiguously
refers to the local declaration for \mbox{\tt sin}.

\subsection{Closure}
\label{closure}
\index{closure}

Every module in a \Haskell{} program must be {\em closed}.  That is,
every name explicitly mentioned by the source code
must be either defined locally or imported from another module.
However, entities that the compiler requires for type checking or other
compile time analysis need not be imported if they are not mentioned
by name.  The \Haskell{} compilation system is responsible for finding
any information needed for compilation without the help of the
programmer.  That is, the import of a variable \mbox{\tt x} does not
require that the datatypes and classes in the signature of \mbox{\tt x} be
brought into the module along with \mbox{\tt x} unless these entities are
referenced by name in the user program.  The \Haskell{}
system silently imports any information that must accompany an
entity for type checking or any other purposes.  Such entities need
not even be explicitly exported: the following program is valid even though
\mbox{\tt T} does not escape \mbox{\tt M1}:
\bprog
{\small\begin{verbatim}
  module M1(x) where
    data T = T
    x = T
  
  module M2 where
    import M1(x)
    y = x
\end{verbatim}}
\eprog
In this example, there is no way to supply an explicit type signature
for \mbox{\tt y} since \mbox{\tt T} is not in scope.
Whether or not \mbox{\tt T} is explicitly exported, module \mbox{\tt M2} knows
enough about \mbox{\tt T} to correctly type check the program.

The type of an exported entity is unaffected by non-exported type
synonyms.  For example, in
\bprog
{\small\begin{verbatim}
  module M(x) where
    type T = Int
    x :: T
    x = 1
\end{verbatim}}
\eprog
the type of \mbox{\tt x} is both \mbox{\tt T} and \mbox{\tt Int}; these are interchangeable even
when \mbox{\tt T} is not in scope.  That is, the definition of \mbox{\tt T} is available
to any module that encounters it whether or not the name \mbox{\tt T} is
in scope.  The only reason to export \mbox{\tt T} is to allow other modules to
refer it by name; the type checker finds the definition of \mbox{\tt T} if
needed whether or not it is exported.

\section{Standard Prelude}
\label{standard-prelude}
\index{standard prelude}
\index{libraries}
Many of the features of \Haskell{} are defined in \Haskell{}
itself as a library of standard datatypes, classes, and
functions, called the ``Standard Prelude.''  In
\Haskell{}, the Prelude is contained in the
module \mbox{\tt Prelude}.\indexmodule{Prelude} There are also
many predefined library modules, which provide less frequently used
functions and types.  For example, complex numbers, arrays, 
and most of the input/output are all part of the standard
libraries.    These are 
defined in Part~\ref{libraries}.
Separating
libraries from the Prelude has the advantage of reducing the size and
complexity of the Prelude, allowing it to be more easily assimilated,
and increasing the space of useful names available to the programmer.

Prelude and library modules differ from other modules in that
their semantics (but not their implementation) are a fixed part of the
\Haskell{} language definition.
This means, for example, that a compiler may optimize calls to
functions in the Prelude without consulting the source code
of the Prelude.

\subsection{The \texorpdfstring{\mbox{\tt Prelude}}{Prelude} Module}
\indexmodule{Prelude}
\index{Prelude!implicit import of}

The \mbox{\tt Prelude} module is imported automatically into all modules as if
by the statement `\mbox{\tt import\ Prelude}', if and only if it is not imported
with an explicit \mbox{\tt import} declaration. This provision for explicit
import allows entities defined in the Prelude to be selectively imported,
just like those from any other module.

The semantics of the entities in \mbox{\tt Prelude} is specified by a reference
implementation of \mbox{\tt Prelude} written in \Haskell{}, given in
Chapter~\ref{stdprelude}.  Some datatypes (such as \mbox{\tt Int}) and
functions (such as \mbox{\tt Int} addition) cannot be specified directly in
\Haskell{}.  Since the treatment of such entities depends on the
implementation, they are not formally defined in Chapter~\ref{stdprelude}.
The implementation of
\mbox{\tt Prelude} is also incomplete in its treatment of tuples: there should
be an infinite family of tuples and their instance declarations, but the
implementation only gives a scheme.

Chapter~\ref{stdprelude} defines the module \mbox{\tt Prelude} using
several other modules: \mbox{\tt PreludeList}, \mbox{\tt PreludeIO}, and so on.
These modules are {\em not} part of Haskell, and they cannot be imported
separately.  They are simply 
there to help explain the structure of the \mbox{\tt Prelude} module; they
should be considered part of its implementation, not part of the language
definition.
 
\subsection{Shadowing Prelude Names}
\label{std-prel-shadowing}

The rules about the Prelude have been cast so that it is
possible to use Prelude names for nonstandard purposes; however,
every module that does so must have an \mbox{\tt import} declaration
that makes this nonstandard usage explicit.  For example:
\bprog
{\small\begin{verbatim}
  module A( null, nonNull ) where
    import Prelude hiding( null ) 
    null, nonNull :: Int -> Bool
    null    x = x == 0
    nonNull x = not (null x)
\end{verbatim}}
\eprog
Module \mbox{\tt A} redefines \mbox{\tt null}, and contains an unqualified reference to \mbox{\tt null}
on the right hand side of \mbox{\tt nonNull}. The latter would be ambiguous
without the \mbox{\tt hiding(null)} on the \mbox{\tt import\ Prelude} statement. Every
module that imports \mbox{\tt A} unqualified, and then makes an unqualified
reference to \mbox{\tt null} must also resolve the ambiguous use of \mbox{\tt null} just as
\mbox{\tt A} does. Thus there is little danger of accidentally shadowing Prelude
names.

It is possible to construct and use a different module to serve in
place of the Prelude.  Other than the fact that it is implicitly
imported, the Prelude is an ordinary \Haskell{} module; it is special
only in that some objects in the Prelude are referenced by special
syntactic constructs.  Redefining names used by the Prelude does not
affect the meaning of these special constructs.  For example, in
\bprog
{\small\begin{verbatim}
  module B where
    import Prelude()
    import MyPrelude
    f x = (x,x)
    g x = (,) x x
    h x = [x] ++ []
\end{verbatim}}
\eprog
the explicit \mbox{\tt import\ Prelude()} declaration prevents the automatic
import of \mbox{\tt Prelude}, while the declaration \mbox{\tt import\ MyPrelude} brings the
non-standard prelude into scope.  
The special syntax for tuples (such as \mbox{\tt (x,x)} and \mbox{\tt (,)}) and lists
(such as \mbox{\tt [x]} and \mbox{\tt []}) continues to refer to the tuples and lists
defined by the standard \mbox{\tt Prelude};
there is no way to redefine the meaning of \mbox{\tt [x]}, for example, in terms of a
different implementation of lists.
On the other hand, the use of \mbox{\tt ++} is not special syntax, so it refers
to \mbox{\tt ++} imported from \mbox{\tt MyPrelude}.

It is not possible, however, to hide \mbox{\tt instance} declarations in the
\mbox{\tt Prelude}.  For example, one cannot define a new instance for \mbox{\tt Show\ Char}.

\section{Separate Compilation}
\index{separate compilation}
Depending on the \Haskell{} implementation used, separate compilation
of mutually recursive modules may require that imported modules contain
additional information so that they may be referenced before they are
compiled.  Explicit type signatures for all exported values may be
necessary to deal with mutual recursion.  The
precise details of separate compilation are not defined by this
report. 

\section{Abstract Datatypes}
\label{abstract-types}

\index{abstract datatype}
The ability to export a datatype without its constructors
allows the construction of abstract datatypes (ADTs).  For example,
an ADT for stacks could be defined as:
\bprog
{\small\begin{verbatim}
  module Stack( StkType, push, pop, empty ) where
    data StkType a = EmptyStk | Stk a (StkType a)
    push x s = Stk x s
    pop (Stk _ s) = s
    empty = EmptyStk
\end{verbatim}}
\eprog
Modules importing \mbox{\tt Stack} cannot construct values of type \mbox{\tt StkType}
because they do not have access to the constructors of the type.
Instead, they must use \mbox{\tt push}, \mbox{\tt pop}, and \mbox{\tt empty} to construct such values.

It is also possible to build an ADT on top of an existing type by
using a \mbox{\tt newtype} declaration.  For example, stacks can be defined
with lists: 
\bprog
{\small\begin{verbatim}
  module Stack( StkType, push, pop, empty ) where
    newtype StkType a = Stk [a]
    push x (Stk s) = Stk (x:s)
    pop (Stk (_:s)) = Stk s
    empty = Stk []
\end{verbatim}}
\eprogNoSkip


%**~footer

% Local Variables: 
% mode: latex
% End:
